\documentclass[preprint,12pt]{elsarticle}

\usepackage{amsmath,amssymb,bm}
\usepackage{graphicx}
\usepackage{booktabs}
\usepackage{siunitx}
\usepackage{hyperref}
\hypersetup{hidelinks}
\usepackage{xcolor}
\usepackage{listings}
\usepackage{lineno}
\usepackage{float}
\usepackage{enumitem}
\usepackage{microtype}
\usepackage{algorithm}
\usepackage{algpseudocode}
\usepackage{caption}
\captionsetup[figure]{skip=4pt}
\journal{Computer Physics Communications}

\lstset{
  basicstyle=\ttfamily\small,
  frame=single,
  breaklines=true,
  columns=fullflexible,
  keepspaces=true
}

\begin{document}

\begin{frontmatter}

\title{A Reproducible Computational Framework for Green-Function and Legendre-Series Solutions in Spherical Electrostatics}

\author[ppu]{Othman H. Y. Zalloum\corref{cor1}}
\ead{ozalloum@ppu.edu}
\cortext[cor1]{Corresponding author.}
\address[ppu]{Department of Applied Mathematics and Physics, Palestine Polytechnic University, Hebron, Palestine}

\begin{abstract}
The electrostatic potential of a point charge inside a grounded conducting sphere is a canonical boundary-value problem with an exact image-charge Green-function solution and an equivalent separated representation in Legendre polynomials. This article presents a reproducible computational framework for evaluating both representations, generating machine-readable data, regenerating publication-quality figures, and validating primary and post-processed electrostatic quantities. The framework evaluates the potential, induced surface charge density, image energy, signed radial force, multipole coefficient decay, computational-cost trends, and spatial truncation-error distributions. In dimensionless units, finite Legendre-series truncations are compared with the exact image-charge solution for several source positions and multipole orders. The relative interior potential error decreases monotonically as the maximum multipole order is increased from $L=2$ to $L=64$ for all numerical cases considered. At $L=64$, the measured relative errors are $1.6556\times10^{-2}$, $3.7953\times10^{-2}$, $4.3527\times10^{-2}$, and $6.5841\times10^{-2}$ for $b/a=0.25$, $0.45$, $0.65$, and $0.80$, respectively. For the induced surface charge at $b/a=0.45$, the relative root-mean-square error decreases from $6.189\times10^{-2}$ at $L=4$ to $3.443\times10^{-11}$ at $L=32$. The package is designed for computational electrostatics, spectral approximation studies, and validation of numerical Poisson/Laplace solvers where exact or semi-analytical reference solutions are required.
\end{abstract}

\begin{keyword}
Green's functions \sep electrostatics \sep image charges \sep grounded conducting sphere \sep Legendre polynomials \sep reproducible computing \sep spectral approximation \sep Poisson equation \sep validation \sep open-source software
\end{keyword}

\end{frontmatter}

\linenumbers

\section*{Program summary}
\noindent\textbf{Program title:} Spherical Green--Legendre Electrostatics Solver\\
\textbf{CPC library link to program files:} To be assigned after acceptance.\\
\textbf{Developer's repository link:} To be supplied after deposition in a public repository.\\
\textbf{Code Ocean capsule:} Not applicable at the present stage.\\
\textbf{Licensing provisions:} MIT, BSD-3-Clause, or GPL-3.0 license recommended for final release.\\
\textbf{Programming language:} Python 3.\\
\textbf{External routines/libraries:} NumPy, SciPy, Matplotlib, pandas.\\
\textbf{Nature of problem:} Accurate evaluation of electrostatic Green-function and Legendre-series solutions for a point charge inside a grounded conducting sphere. The computation involves singular-field evaluation, boundary enforcement on a curved conducting surface, modal truncation, induced surface-charge recovery, energy and force post-processing, and multipole-coefficient decay.\\
\textbf{Solution method:} The exact image-charge solution is evaluated directly and compared with finite Legendre-series truncations. The same reproducibility datasets are used for figures, tables, and validation diagnostics. Interior field errors, boundary-derived surface-charge errors, image energy, signed radial force, coefficient decay, computational cost, and spatial error maps are generated by a single scripted workflow.\\
\textbf{Restrictions:} The current implementation treats an axisymmetric point charge located on the symmetry axis of a perfectly conducting grounded sphere. Extensions to off-axis sources, dielectric interfaces, and non-spherical geometries are natural future developments.\\
\textbf{Unusual features:} The workflow evaluates both primary fields and derived physical observables using analytically controlled reference solutions, while also exporting publication-quality figures and machine-readable data files.

\section{Introduction}

Green-function methods and image-charge constructions are central tools in classical electrostatics. They provide exact solutions for selected boundary-value problems and give direct physical insight into how conducting boundaries modify electric fields, surface charge distributions, energies, and forces \cite{jackson1999,zangwill2013,reitz2008,stratton1941,landau1984}. Recent work has continued to develop Green-function and image-method viewpoints for spherical and spheroidal conducting geometries, confirming the continued relevance of analytically controlled electrostatic solutions in modern boundary-value analysis \cite{deoliveira2012,goldstein2018,vucic2026}. Such analytically controlled configurations are also valuable in computational physics because they provide reference solutions against which numerical approximations can be checked without ambiguity.

A point charge inside a grounded conducting sphere is a standard example of this type. The problem is simple enough to admit a closed image-charge solution, but it is also rich enough to test several numerical features. The potential contains a point singularity, the grounded boundary condition must be enforced on a curved surface, and physically important derived quantities such as induced surface charge density and image force depend on derivatives or limiting values of the field. These features make the problem more informative than a purely smooth test case.

The same physical solution can also be obtained by separation of variables. For an on-axis charge, the potential is represented by an infinite Legendre-polynomial expansion. Truncating this expansion at finite multipole degree $L$ gives a controlled spectral approximation. The comparison between the exact Green-function solution and the truncated Legendre-series solution therefore gives a direct way to quantify spectral truncation error, boundary-derived errors, and multipole-coefficient decay. Such questions are relevant to computational electromagnetics, spectral methods, boundary-element and finite-element formulations, and numerical solutions of Poisson and Laplace equations \cite{beer2008,liu2010,jin2014,cheng2005,canuto2007,boyd2001,gottlieb1977,trefethen2000,trefethen2019,pittino2014}.

This paper reformulates the classical grounded-sphere problem as a reproducible computational framework suitable for computational physics software validation. The contribution is not the image solution itself, which is classical, but the integration of exact-field evaluation, finite-order spectral approximation, error diagnostics, reproducibility data, computational-cost reporting, spatial error mapping, and physical post-processing into a compact open-source workflow. The implementation is intended to support users who need transparent reference solutions for electrostatic Poisson/Laplace solvers and for spectral expansions in spherical geometry. The use of exact or semi-analytical solutions as reference fields is a common strategy in computational physics, particularly when numerical schemes must be assessed for boundary enforcement, long-range interactions, derivative recovery, and post-processed physical observables \cite{mazars2011,westerberg2023,dantchev2023}.

The source position plays an important role. Legendre expansions are widely used in spherical and nearly spherical geometries because they expose the modal content of the solution and the rate at which higher-order angular structure is suppressed \cite{arfken2013,boyd1992}. In the Legendre representation, the regular harmonic correction contains coefficients proportional to $b^\ell$, where $b/a$ is the dimensionless source radius. As the charge moves closer to the conducting boundary, high-order multipoles decay more slowly and the induced surface charge becomes increasingly localized near the closest point on the surface. This behavior provides a demanding validation case for finite-order approximations and for numerical solvers that attempt to recover both the field and its boundary derivative.

The remainder of the paper is organized as follows. Section~\ref{sec:model} defines the electrostatic model problem and gives the exact image-charge Green-function solution. Section~\ref{sec:legendre} gives the Legendre-series representation and its finite-$L$ truncation. Section~\ref{sec:observables} summarizes the induced surface charge, image energy, and force formulas. Section~\ref{sec:architecture} describes the software architecture and data organization. Section~\ref{sec:implementation} describes the computational implementation. Section~\ref{sec:errors} defines the validation methodology and error measures. Section~\ref{sec:results} presents numerical results and the complete figure set. Section~\ref{sec:performance} discusses computational cost and reproducibility. Section~\ref{sec:discussion} discusses implications for computational electrostatics and spectral approximation.

\section{Model problem and exact Green-function solution}\label{sec:model}

Consider a grounded conducting sphere of radius $a$. A point charge $q$ is located at
\begin{equation}
    \mathbf r_0=b\hat{\mathbf z}, \qquad 0<b<a .
\end{equation}
The potential inside the sphere satisfies
\begin{equation}
    \nabla^2\Phi(\mathbf r)=-\frac{q}{\epsilon_0}\delta(\mathbf r-\mathbf r_0),
    \qquad \Phi(a,\theta)=0 .
    \label{eq:poisson}
\end{equation}
Throughout the numerical part of this paper we use dimensionless units
\begin{equation}
    a=1, \qquad q=1, \qquad \frac{1}{4\pi\epsilon_0}=1 .
\end{equation}
Dimensional values may be restored by multiplying the dimensionless potential by $q/(4\pi\epsilon_0 a)$.

The grounded sphere image construction, also discussed in modern treatments of related grounded-sphere configurations \cite{deoliveira2012}, places an image charge
\begin{equation}
    q'=-\frac{a}{b}q
\end{equation}
at the exterior point $b'=a^2/b$. In the normalized units used below, the exact potential is
\begin{equation}
\Phi_{\rm img}(r,\theta)=
\frac{1}{\sqrt{r^2+b^2-2rb\cos\theta}}
-
\frac{1/b}{\sqrt{r^2+b^{-2}-2rb^{-1}\cos\theta}} .
\label{eq:image}
\end{equation}
This expression satisfies the grounded boundary condition at $r=1$ and has the required Coulomb singularity at the physical charge position.

\section{Legendre-series representation}\label{sec:legendre}

For an on-axis charge, the solution is axisymmetric and can be represented with Legendre polynomials. The Coulomb kernel expansion is
\begin{equation}
\frac{1}{|\mathbf r-\mathbf r_0|}=\sum_{\ell=0}^{\infty}
\frac{r_<^\ell}{r_>^{\ell+1}}P_\ell(\cos\theta),
\end{equation}
where $r_<=\min(r,b)$ and $r_>=\max(r,b)$. Adding the regular harmonic correction required to enforce the grounded boundary gives
\begin{equation}
\Phi_L(r,\theta)=\sum_{\ell=0}^{L}\left[
\frac{r_<^\ell}{r_>^{\ell+1}}-b^\ell r^\ell
\right]P_\ell(\cos\theta).
\label{eq:legendre_trunc}
\end{equation}
The limit $L\rightarrow\infty$ gives the exact solution away from the point-charge singularity. The finite-$L$ approximation is the numerical object tested in the reproducibility data. Standard properties of Legendre polynomials and special-function expansions are documented in classical references and modern mathematical-physics texts \cite{abramowitz1972,gradshteyn2007,szego1939,arfken2013}.

\section{Induced surface charge, energy, and force}\label{sec:observables}

The induced surface charge density is obtained from the outward normal derivative of the potential at the conductor:
\begin{equation}
    \sigma(\theta)=-\epsilon_0\left.\frac{\partial\Phi}{\partial r}\right|_{r=a} .
\end{equation}
In the normalized units used for the numerical tables, the exact surface charge density is
\begin{equation}
\sigma(\theta)=-\frac{1-b^2}{4\pi\left(1+b^2-2b\cos\theta\right)^{3/2}} .
\label{eq:sigma_exact}
\end{equation}
It satisfies the charge-sum rule
\begin{equation}
    2\pi\int_0^\pi \sigma(\theta)\sin\theta\,d\theta=-1 .
\label{eq:charge_sum}
\end{equation}
The image contribution to the potential at the charge position gives the image energy
\begin{equation}
    U(b)=-\frac{1}{2(1-b^2)} .
\label{eq:energy}
\end{equation}
The numerical energy--force data use the signed attraction convention
\begin{equation}
    F_{\rm data}(b)=\frac{dU}{db}=-\frac{b}{(1-b^2)^2},
\label{eq:force}
\end{equation}
so negative values denote attraction toward the conducting boundary in the convention used for the force plot. The physical outward radial force magnitude is $|F_{\rm data}|$.


\section{Software architecture and data organization}\label{sec:architecture}

The framework is organized as a small, inspectable computational physics package rather than as a collection of independent plotting scripts. This organization follows the practical need, common in computational electromagnetics and boundary-value solvers, to separate model equations, numerical evaluation, diagnostics, and reproducibility artifacts \cite{cheng2005,pittino2014}. The design separates analytical formula evaluation, finite-order spectral evaluation, diagnostic computation, data export, figure generation, and metadata recording. This organization is intended to make the numerical results auditable and to allow individual components to be reused in external electrostatics codes.

The recommended package structure is
\begin{lstlisting}[language=bash,basicstyle=\ttfamily\small,frame=single]
spherical-green-legendre/
  README.md
  requirements.txt
  environment.yml
  src/
    image_solution.py
    legendre_solution.py
    observables.py
    errors.py
  data/
    figure_01_data.csv
    ...
  figures/
    figure_01.pdf
    ...
  scripts/
    generate_all_results_cpc.py
  tests/
    test_boundary_condition.py
    test_charge_conservation.py
    test_energy_force.py
  metadata/
    run_summary.json
\end{lstlisting}
The source-code directory contains the numerical routines. The data directory contains the machine-readable numerical values used in the article. The figure directory contains vector PDF figures and preview PNG files. The script directory contains the end-to-end reproduction script. The metadata directory records numerical settings such as source positions, truncation orders, grid definitions, and library versions.

The package exposes four conceptual layers. The first layer evaluates exact analytical expressions. The second layer evaluates finite Legendre-series approximations. The third layer computes physical observables and numerical diagnostics. The fourth layer exports reproducibility data and publication figures. This separation is important because users may wish to compare their own numerical Poisson/Laplace solver only against the exact potential, only against the surface charge, or against the full set of diagnostics.


\section{Computational implementation}\label{sec:implementation}

The computational workflow is designed so that every numerical value in the Results section can be regenerated from explicit input parameters. The implementation consists of the following components:
\begin{enumerate}[label=(\roman*)]
    \item direct evaluation of Eq.~\eqref{eq:image} on Cartesian and spherical grids;
    \item finite-$L$ evaluation of Eq.~\eqref{eq:legendre_trunc};
    \item masking of the point-charge singularity when computing field-norm errors;
    \item evaluation of the induced surface charge from Eq.~\eqref{eq:sigma_exact} and from differentiated truncated series;
    \item generation of reproducibility datasets for field plots, line profiles, convergence curves, surface-charge curves, energy/force curves, multipole coefficients, computational cost, and spatial error maps;
    \item automated regeneration of all figures and tables from the stored datasets.
\end{enumerate}

The public package is organized into source-code, data, figure, script, test, and metadata directories. Figure~\ref{fig:workflow} summarizes the computational workflow and Fig.~\ref{fig:repro_structure} summarizes the intended reproducibility structure.

\begin{algorithm}[t]
\caption{End-to-end reproduction workflow.}\label{alg:workflow}
\begin{algorithmic}[1]
\State Define source positions $b/a$, truncation orders $L$, angular grid, radial grid, and singular-mask parameters.
\For{each source position $b/a$}
    \State Evaluate the image-charge reference solution $\Phi_{\rm img}$.
    \For{each truncation order $L$}
        \State Evaluate the finite Legendre solution $\Phi_L$.
        \State Compute masked field errors and export convergence data.
    \EndFor
    \State Evaluate exact and truncated induced surface-charge densities.
    \State Compute image energy, signed radial force, and multipole coefficient decay.
\EndFor
\State Generate computational-cost data and spatial error maps.
\State Export CSV/JSON datasets, metadata, and vector figures.
\end{algorithmic}
\end{algorithm}

\begin{lstlisting}[language=Python,caption={Minimal pseudocode representation of the solver loop.},label={lst:pseudocode}]
for b in source_positions:
    phi_exact = image_solution(grid, b)
    for L in truncation_orders:
        phi_L = legendre_solution(grid, b, L)
        E_L = relative_l2_error(phi_L, phi_exact, mask)
        save_error(b, L, E_L)

sigma_exact = exact_surface_charge(theta, b)
for L in surface_charge_orders:
    sigma_L = truncated_surface_charge(theta, b, L)
    save_surface_error(L, sigma_L, sigma_exact)
\end{lstlisting}

The direct image solution is evaluated from Eq.~\eqref{eq:image}. The Legendre solution is evaluated from Eq.~\eqref{eq:legendre_trunc} by summing terms from $\ell=0$ to $L$. The surface-charge calculation uses the analytical surface expression and the differentiated finite series. The package therefore tests both field evaluation and derivative-based post-processing.

\section{Validation methodology and error measures}\label{sec:errors}

The numerical study is defined by reproducibility datasets corresponding to the reported figures and tables. These datasets contain plotted field values, one-dimensional line profiles, relative potential errors versus truncation order, exact and truncated surface charge densities, surface-charge distributions for several source positions, image energy and signed radial force values, Legendre coefficient decay, computational-cost estimates, and a spatial error map. The same datasets are used for figures and numerical values.

The principal field-error diagnostic is the relative potential error
\begin{equation}
    E_L=\frac{\|\Phi_L-\Phi_{\rm img}\|_2}{\|\Phi_{\rm img}\|_2} .
\label{eq:error}
\end{equation}
The quadrature and masking procedure are those used to generate the tabulated convergence data. For final repository deposition, the grid size, quadrature rule, and mask radius should be recorded in machine-readable metadata.

For boundary accuracy, the exact induced surface charge density is compared with the surface charge obtained from the differentiated truncated series. Similar attention to derived quantities is important in computational electrostatics because physically relevant outputs often depend on derivatives, surface fields, or integral constraints rather than only on the scalar potential itself \cite{jin2014,liu2026}. This diagnostic is included because boundary derivatives often converge differently from the field itself and can reveal truncation effects that are not visible in a potential plot. Additional diagnostics are provided by the image energy, the signed radial force, and the coefficient relation
\begin{equation}
    |c_\ell|=|b|^\ell .
\end{equation}

\section{Results}\label{sec:results}

\subsection{Primary field and spectral convergence}

Figure~\ref{fig:field} shows the exact image-charge potential in the $x$--$z$ meridional plane. Values outside the sphere are blank in the tabulated dataset. The plotted field is restricted to the physical domain $r/a\leq1$, and the potential is clipped for visualization near the point-charge singularity.

\begin{figure}[t]
\centering
\includegraphics[width=0.68\linewidth]{figures/figure_01.pdf}
\caption{Exact image-charge potential in the meridional plane of the grounded conducting sphere. The reproducibility dataset contains $x/a$, $z/a$, $r/a$, and a clipped potential value; entries outside the sphere are blank.}
\label{fig:field}
\end{figure}

\begin{figure}[t]
\centering
\includegraphics[width=0.68\linewidth]{figures/figure_02.pdf}
\caption{Exact and truncated Legendre-series potentials along the line used in the reproducibility data. The plotted truncation orders are $L=2$, $4$, $8$, $16$, and $32$.}
\label{fig:line_profile}
\end{figure}

Figure~\ref{fig:line_profile} compares the exact potential with truncation orders $L=2$, $4$, $8$, $16$, and $32$. The measured relative RMS errors from this one-dimensional dataset decrease monotonically with $L$, as shown in Table~\ref{tab:line_errors}.

\begin{table}[t]
\centering
\caption{Errors in the line-comparison data. The relative RMS value is computed from the corresponding tabulated line-profile data.}
\label{tab:line_errors}
\begin{tabular}{ccc}
\toprule
$L$ & Relative RMS error & Maximum absolute error \\
\midrule
2  & $2.346\times10^{-1}$ & $4.565\times10^{-1}$ \\
4  & $1.553\times10^{-1}$ & $3.690\times10^{-1}$ \\
8  & $9.155\times10^{-2}$ & $2.807\times10^{-1}$ \\
16 & $5.003\times10^{-2}$ & $2.039\times10^{-1}$ \\
32 & $2.616\times10^{-2}$ & $1.426\times10^{-1}$ \\
\bottomrule
\end{tabular}
\end{table}

\begin{figure}[t]
\centering
\includegraphics[width=0.68\linewidth]{figures/figure_03.pdf}
\caption{Relative $L^2$ error of the truncated Legendre expansion as a function of maximum multipole degree $L$ for $b/a=0.25$, $0.45$, $0.65$, and $0.80$. The data show monotonic error reduction from $L=2$ to $L=64$ for all sampled source positions.}
\label{fig:l2_error}
\end{figure}

The relative $L^2$ errors in Fig.~\ref{fig:l2_error} are available for $b/a=0.25$, $0.45$, $0.65$, and $0.80$, with $L=2$, $4$, $8$, $16$, $32$, and $64$. For all four source positions, the error decreases monotonically over the sampled range of $L$. The error at $L=64$ remains finite in the tabulated data and is not at machine precision. Therefore, no claim of $10^{-6}$ convergence is made here.

\begin{table}[t]
\centering
\caption{Measured relative $L^2$ errors from the corresponding tabulated convergence data.}
\label{tab:l2_errors}
\begin{tabular}{cccc}
\toprule
$b/a$ & $E_2$ & $E_{64}$ & Reduction factor $E_2/E_{64}$ \\
\midrule
0.25 & 0.362593 & 0.016556 & 21.90 \\
0.45 & 0.551051 & 0.037953 & 14.52 \\
0.65 & 0.699394 & 0.043527 & 16.07 \\
0.80 & 0.815765 & 0.065841 & 12.39 \\
\bottomrule
\end{tabular}
\end{table}

\subsection{Boundary observables and physical diagnostics}

\begin{figure}[t]
\centering
\includegraphics[width=0.68\linewidth]{figures/figure_04.pdf}
\caption{Exact and truncated induced surface charge density for the sampled case, plotted using the same signed convention. The truncation orders are $L=4$, $8$, $16$, and $32$.}
\label{fig:sigma_trunc}
\end{figure}

Figure~\ref{fig:sigma_trunc} compares the exact induced surface charge density with truncated approximations. In this dataset the convergence is rapid for $b/a=0.45$: the relative RMS error decreases from $6.189\times10^{-2}$ at $L=4$ to $3.443\times10^{-11}$ at $L=32$. The maximum absolute errors are reported in Table~\ref{tab:sigma_errors}.

\begin{table}[t]
\centering
\caption{Surface-charge errors computed against the exact surface-charge data.}
\label{tab:sigma_errors}
\begin{tabular}{ccc}
\toprule
$L$ & Maximum absolute error & Relative RMS error \\
\midrule
4  & $3.374\times10^{-2}$  & $6.189\times10^{-2}$ \\
8  & $2.259\times10^{-3}$  & $3.471\times10^{-3}$ \\
16 & $6.745\times10^{-6}$  & $8.321\times10^{-6}$ \\
32 & $3.573\times10^{-11}$ & $3.443\times10^{-11}$ \\
\bottomrule
\end{tabular}
\end{table}

\begin{figure}[t]
\centering
\includegraphics[width=0.68\linewidth]{figures/figure_05.pdf}
\caption{Exact induced surface charge density for $b/a=0.25$, $0.45$, $0.65$, and $0.80$. The peak magnitude at $\theta=0$ increases as the charge moves toward the conducting boundary.}
\label{fig:sigma_positions}
\end{figure}

\begin{table}[t]
\centering
\caption{Surface-charge diagnostics from the corresponding tabulated surface-charge data.}
\label{tab:sigma_diag}
\begin{tabular}{cccc}
\toprule
$b/a$ & $\sigma(0)$ & $\sigma(\pi)$ & Numerical total charge \\
\midrule
0.25 & -0.176839 & -0.038197 & -0.999997 \\
0.45 & -0.381446 & -0.020817 & -0.999994 \\
0.65 & -1.071860 & -0.010230 & -0.999984 \\
0.80 & -3.580986 & -0.004912 & -0.999962 \\
\bottomrule
\end{tabular}
\end{table}

Figure~\ref{fig:sigma_positions} shows exact induced surface charge densities for $b/a=0.25$, $0.45$, $0.65$, and $0.80$. The density is most negative at $\theta=0$ and least negative at $\theta=\pi$. The trapezoidal integrations of the tabulated data give total induced charges close to $-1$, as expected from the grounded-conductor solution.

\begin{figure}[t]
\centering
\includegraphics[width=0.68\linewidth]{figures/figure_06.pdf}
\caption{Image energy $U(b)=-1/[2(1-b^2)]$ and signed radial force $F_{\rm data}(b)=-b/(1-b^2)^2$ over the sampled interval $0.02\leq b/a\leq0.95$.}
\label{fig:energy_force}
\end{figure}

Figure~\ref{fig:energy_force} reports the image energy and the signed radial force convention in Eq.~\eqref{eq:force} for $0.02\leq b/a\leq0.95$. Over this interval, the tabulated image energy decreases from $-0.500200$ to $-5.128205$, while the signed force decreases from $-0.020016$ to $-99.934254$. These values match Eqs.~\eqref{eq:energy} and \eqref{eq:force} evaluated at the tabulated endpoints.

\begin{figure}[t]
\centering
\includegraphics[width=0.68\linewidth]{figures/figure_07.pdf}
\caption{Magnitude of the Legendre multipole coefficients $|b|^\ell$ for $b/a=0.25$, $0.45$, $0.65$, and $0.80$, sampled from $\ell=0$ to $80$.}
\label{fig:coeff_decay}
\end{figure}

Figure~\ref{fig:coeff_decay} shows $|b|^\ell$ for $b/a=0.25$, $0.45$, $0.65$, and $0.80$ over $0\leq\ell\leq80$. At $\ell=80$, the tabulated coefficients are $6.842\times10^{-49}$, $1.807\times10^{-28}$, $1.079\times10^{-15}$, and $1.767\times10^{-8}$, respectively. The coefficient decay is therefore much slower for $b/a=0.80$ than for charges closer to the center.

\subsection{Computational workflow, cost, and spatial error}

\begin{figure}[t]
\centering
\includegraphics[width=0.95\linewidth]{figures/figure_08_computational_workflow.pdf}
\caption{Computational workflow used to generate the reproducibility datasets and publication figures. The workflow begins with input parameters, evaluates the exact image solution and the finite Legendre representation, computes physical observables and error diagnostics, and exports machine-readable data tables and publication figures.}
\label{fig:workflow}
\end{figure}

Figure~\ref{fig:workflow} emphasizes that the study is organized as a reproducible software workflow rather than a single set of manually produced plots. The same input parameters drive the exact-solution evaluation, Legendre-series evaluation, field diagnostics, boundary diagnostics, and output generation.

\begin{figure}[t]
\centering
\includegraphics[width=0.72\linewidth]{figures/figure_09_computational_cost.pdf}
\caption{Representative computational cost as a function of maximum multipole degree $L$. The data quantify the cost of increasing the truncation order in the Legendre-series solver.}
\label{fig:cost}
\end{figure}

Figure~\ref{fig:cost} reports the computational cost trend associated with increasing $L$. This is useful for users who need to balance accuracy against runtime when selecting the truncation order.

\begin{figure}[t]
\centering
\includegraphics[width=0.72\linewidth]{figures/figure_10_error_heatmap.pdf}
\caption{Spatial distribution of the absolute truncation error $|\Phi_{\rm img}-\Phi_L|$ for the selected source position and truncation order. The map identifies where the finite Legendre representation differs most strongly from the image-charge reference solution.}
\label{fig:error_heatmap}
\end{figure}

Figure~\ref{fig:error_heatmap} complements the scalar error curves by showing where the truncation error is concentrated inside the physical domain. Such maps are helpful for diagnosing whether an error metric is dominated by the singular neighborhood, by near-boundary behavior, or by broad interior differences.

\begin{figure}[t]
\centering
\includegraphics[width=0.95\linewidth]{figures/figure_11_reproducibility_structure.pdf}
\caption{Recommended reproducibility structure for the public package. The organization separates the notebook, figures, data, tables, generated outputs, package instructions, dependency files, and the regeneration script so that the numerical results can be rerun and audited.}
\label{fig:repro_structure}
\end{figure}

Figure~\ref{fig:repro_structure} summarizes the intended repository organization. This structure supports independent reruns, package inspection, and long-term reuse of the data and scripts.


\section{Performance analysis and reproducibility}\label{sec:performance}

The computational cost is dominated by the finite Legendre summation over the selected grid. For the sampled implementation and grid, the runtime increases monotonically with the maximum multipole degree $L$; the measured cost curve is reported directly rather than fitted to a presumed asymptotic scaling law. The representative cost curve in Fig.~\ref{fig:cost} therefore provides a practical guide for selecting $L$ in reproducible calculations. The exact image solution is inexpensive by comparison, but it remains essential because it supplies the reference field used by all error metrics.

The spatial error map in Fig.~\ref{fig:error_heatmap} is included to complement scalar error norms. A single global norm can hide the location of the largest discrepancies, particularly when a singular source and a curved boundary are present. The heat map identifies whether the finite-order representation differs most strongly near the source, near the conductor, or across a broader interior region. This information is useful when adapting the framework to finite-element, finite-difference, boundary-element, or spectral solvers.

Reproducibility is supported by three design choices. First, the package exports the numerical datasets used to generate each figure. Second, a single script regenerates the complete figure set from the stored parameters. Third, metadata files record the numerical configuration so that later users can distinguish scientific changes from changes in plotting, library versions, or numerical resolution. For final submission, the public repository should include a release tag and a persistent archive DOI.

\section{Discussion}\label{sec:discussion}

The results demonstrate how a classical electrostatic boundary-value problem can serve as a compact validation and reproducibility framework for computational physics. The exact image-charge solution provides a closed reference field satisfying the grounded conducting boundary condition, while the finite Legendre-series solution provides a transparent spectral approximation whose truncation error can be measured directly. This separation between exact physics and finite-order representation is useful for validating numerical implementations of Poisson and Laplace solvers, including finite-element, boundary-element, finite-volume, and spectral implementations \cite{beer2008,liu2010,jin2014,cheng2005,pittino2014}.

The potential-error data show monotonic improvement with increasing truncation order for all sampled source positions. At $L=64$, the largest measured relative $L^2$ error occurs for $b/a=0.80$ and is $6.5841\times10^{-2}$, while the smallest occurs for $b/a=0.25$ and is $1.6556\times10^{-2}$. These values show clear convergence over the sampled range, while also indicating that the tested interior-error diagnostic is not at machine precision. This distinction is important: the framework reports the measured numerical behavior rather than assuming an asymptotic rate beyond the sampled data.

The dependence on source position has a direct physical and mathematical interpretation. As the charge approaches the conducting surface, the induced charge density becomes more strongly localized near $\theta=0$, and the multipole coefficients decay more slowly. The coefficient data make this effect explicit: at $\ell=80$, the coefficient for $b/a=0.80$ remains many orders of magnitude larger than the corresponding coefficients for sources closer to the center. Near-boundary sources therefore provide more demanding validation cases for finite-order spectral approximations and for numerical solvers resolving localized boundary response.

The induced surface charge density is a particularly important diagnostic because it is obtained from the normal derivative of the potential at the conducting boundary. A numerical method may reproduce the potential with acceptable accuracy while still producing larger errors in derivative-based quantities. For the sampled case $b/a=0.45$, the differentiated truncated series reproduces the exact surface charge very accurately by $L=32$, with a relative RMS error of $3.443\times10^{-11}$. This result should be interpreted as a measured property of the present dataset and sampling, not as a universal convergence rate for all source positions or all derivative quantities.

The energy, signed radial force, and multipole-coefficient datasets provide independent physical consistency checks. The image energy and signed force agree with the closed-form expressions derived from the Green-function solution, while the coefficient table agrees with the exact relation $|c_\ell|=|b|^\ell$. These diagnostics test different parts of the calculation: field evaluation, boundary differentiation, energetic post-processing, force evaluation, and modal decay. A solver that passes all of these checks provides stronger evidence of correctness than one tested only against the potential field.

The additional computational figures strengthen the connection to \textit{Computer Physics Communications}. The workflow diagram clarifies how the solver generates data, diagnostics, and figures; the cost plot reports the effect of increasing $L$; the error heat map localizes truncation error spatially; and the reproducibility-structure diagram documents how the package should be organized for public release. Together, these additions help position the work as a reusable computational physics resource.

Beyond the grounded sphere itself, the framework is useful for validating computational electrostatic solvers and related numerical methods for elliptic boundary-value problems. The same design philosophy could be extended to other geometries for which Green functions or semi-analytical reference solutions are available, including spherical-surface problems, spheroidal conducting boundaries, and environment-dependent Green-tensor formulations \cite{goldstein2018,vucic2026,westerberg2023}. Boundary-element and finite-element implementations can compare their numerical fields against the exact Green-function solution, while spectral implementations can compare modal truncation behavior directly with the Legendre-series data. Extensions to dielectric spheres, off-axis charges, or nonseparable geometries could provide additional validation cases in which exact or semi-analytical reference solutions remain available.

\section{Conclusion}

This paper presented a reproducible computational framework for a classical electrostatic boundary-value problem: the potential of a point charge inside a grounded conducting sphere. The exact image-charge Green-function solution was used as a reference, and finite Legendre-series truncations were tested through field errors, line profiles, induced surface charge, image energy, signed radial force, multipole-coefficient decay, computational cost, and spatial error mapping.

The tabulated data show monotonic improvement of the relative interior potential error from $L=2$ to $L=64$ for $b/a=0.25$, $0.45$, $0.65$, and $0.80$. For $b/a=0.45$, the induced surface-charge diagnostic converges rapidly over the sampled range from $L=4$ to $L=32$. The image energy, signed force, and modal coefficients agree with their closed-form expressions, providing additional physical validation checks.

The resulting framework provides a compact reference-solution package for validating computational electrostatics, spectral approximations, and numerical solvers for Poisson and Laplace boundary-value problems. Because the reference solution and several derived observables are known analytically, the package can be used to assess both primary-field accuracy and the reliability of post-processed physical quantities in future computational studies.

\section*{Data and code availability}
The numerical values used in the Results section are contained in the reproducibility datasets associated with the figures and tables. For final submission, the author should deposit the source code, datasets, figure scripts, tests, metadata files, and environment file in a public repository and include the repository DOI or persistent URL here.

\section*{Acknowledgements}
The author acknowledges Palestine Polytechnic University for its support.

\section*{Declaration of competing interest}
The author declares that there are no known competing financial interests or personal relationships that could have appeared to influence the work reported in this paper.

\begin{thebibliography}{99}

\bibitem{jackson1999}
J. D. Jackson, \textit{Classical Electrodynamics}, 3rd ed., John Wiley \& Sons, New York, 1999.

\bibitem{zangwill2013}
A. Zangwill, \textit{Modern Electrodynamics}, Cambridge University Press, Cambridge, 2013.

\bibitem{reitz2008}
J. R. Reitz, F. J. Milford, R. W. Christy, \textit{Foundations of Electromagnetic Theory}, 4th ed., Addison-Wesley, Boston, 2008.

\bibitem{stratton1941}
J. A. Stratton, \textit{Electromagnetic Theory}, McGraw-Hill, New York, 1941.

\bibitem{morse1953}
P. M. Morse, H. Feshbach, \textit{Methods of Theoretical Physics}, McGraw-Hill, New York, 1953.

\bibitem{arfken2013}
G. B. Arfken, H. J. Weber, F. E. Harris, \textit{Mathematical Methods for Physicists: A Comprehensive Guide}, 7th ed., Academic Press, 2013.

\bibitem{beer2008}
G. Beer, I. Smith, C. Duenser, \textit{The Boundary Element Method with Programming}, Springer, Vienna, 2008. doi:10.1007/978-3-211-71576-5.

\bibitem{liu2010}
Y. Liu, \textit{Fast Multipole Boundary Element Method: Theory and Applications in Engineering}, Cambridge University Press, Cambridge, 2010. doi:10.1017/CBO9780511605345.

\bibitem{jin2014}
J.-M. Jin, \textit{The Finite Element Method in Electromagnetics}, 3rd ed., Wiley-IEEE Press, Hoboken, 2014.

\bibitem{canuto2007}
C. Canuto, M. Y. Hussaini, A. Quarteroni, T. A. Zang, \textit{Spectral Methods: Fundamentals in Single Domains}, Springer, Berlin, 2007. doi:10.1007/978-3-540-30726-6.

\bibitem{boyd2001}
J. P. Boyd, \textit{Chebyshev and Fourier Spectral Methods}, 2nd ed., Dover Publications, Mineola, 2001.

\bibitem{gottlieb1977}
D. Gottlieb, S. A. Orszag, \textit{Numerical Analysis of Spectral Methods: Theory and Applications}, SIAM, Philadelphia, 1977. doi:10.1137/1.9781611970425.

\bibitem{bernardi1997}
C. Bernardi, Y. Maday, Spectral methods, in: P. G. Ciarlet, J.-L. Lions (Eds.), \textit{Handbook of Numerical Analysis}, vol. 5, North-Holland, Amsterdam, 1997, pp. 209--485.

\bibitem{trefethen2000}
L. N. Trefethen, \textit{Spectral Methods in MATLAB}, SIAM, Philadelphia, 2000. doi:10.1137/1.9780898719598.

\bibitem{boyd1992}
J. P. Boyd, The rate of convergence of Legendre series, \textit{Journal of Computational Physics} 101 (1992) 366--377. doi:10.1016/0021-9991(92)90041-N.

\bibitem{trefethen2019}
L. N. Trefethen, \textit{Approximation Theory and Approximation Practice}, extended ed., SIAM, Philadelphia, 2019. doi:10.1137/1.9781611975949.

\bibitem{gradshteyn2007}
I. S. Gradshteyn, I. M. Ryzhik, \textit{Table of Integrals, Series, and Products}, 7th ed., Academic Press, Boston, 2007.

\bibitem{abramowitz1972}
M. Abramowitz, I. A. Stegun, \textit{Handbook of Mathematical Functions}, Dover Publications, New York, 1972.

\bibitem{szego1939}
G. Szeg{\H o}, \textit{Orthogonal Polynomials}, American Mathematical Society Colloquium Publications, vol. 23, American Mathematical Society, Providence, 1939.

\bibitem{landau1984}
L. D. Landau, E. M. Lifshitz, L. P. Pitaevskii, \textit{Electrodynamics of Continuous Media}, 2nd ed., Butterworth-Heinemann, Oxford, 1984.

\bibitem{goldstein2018}
A. Goldstein, O. Eyal, Relationships among charges, fields, and potential on spherical surfaces boundary value problems, \textit{Results in Physics} 11 (2018) 335--349. doi:10.1016/j.rinp.2018.09.015.

\bibitem{westerberg2023}
N. Westerberg, R. Bennett, Perturbative light--matter interactions: from first principles to inverse design, \textit{Physics Reports} 1026 (2023) 1--63. doi:10.1016/j.physrep.2023.07.005.

\bibitem{deoliveira2012}
M. T. de Oliveira, The method of images applied to the grounded sphere: The problem of the ground wire, \textit{Journal of Electrostatics} 70 (2012) 292--299. doi:10.1016/j.elstat.2012.03.008.

\bibitem{vucic2026}
F. Vu\v{c}i\'c, B. Trkulja, Electrostatic Dirichlet Green function for the exterior boundary value problem outside an oblate spheroid and conducting disc, \textit{Journal of Electrostatics} 141 (2026) 104304. doi:10.1016/j.elstat.2026.104304.

\bibitem{cheng2005}
A. H.-D. Cheng, D. T. Cheng, Heritage and early history of the boundary element method, \textit{Engineering Analysis with Boundary Elements} 29 (2005) 268--302. doi:10.1016/j.enganabound.2004.12.001.

\bibitem{pittino2014}
F. Pittino, L. Selmi, Use and comparative assessment of the CVFEM method for Poisson--Boltzmann and Poisson--Nernst--Planck three dimensional simulations of impedimetric nano-biosensors operated in the DC and AC small signal regimes, \textit{Computer Methods in Applied Mechanics and Engineering} 278 (2014) 902--923. doi:10.1016/j.cma.2014.06.006.

\bibitem{mazars2011}
M. Mazars, Long ranged interactions in computer simulations and for quasi-2D systems, \textit{Physics Reports} 500 (2011) 43--116. doi:10.1016/j.physrep.2010.11.004.

\bibitem{dantchev2023}
D. M. Dantchev, S. Dietrich, Critical Casimir effect: Exact results, \textit{Physics Reports} 1005 (2023) 1--130. doi:10.1016/j.physrep.2022.12.004.

\bibitem{liu2026}
X. Liu, Q. Wu, Y. Jin, F. Hui, H. Zhang, W. Liu, Modeling, analysis, and experimental validation of electromagnetic detumbling for tumbling space targets, \textit{Results in Engineering} 29 (2026) 109630. doi:10.1016/j.rineng.2026.109630.

\end{thebibliography}

\end{document}
